% Paper SCC — Abstract & Structure
% Target: SAFECOMP 2027 (deadline ~Março 2027) or DASC 2027 (deadline ~Junho 2027)

\documentclass[conference]{IEEEtran}
\usepackage{cite,amsmath,amssymb,graphicx,hyperref,listings}

\title{SCC: A Safety Case Compiler for Autonomous Drones}
\author{Andre Oliveira\\RNFC Project\\\texttt{handre.oliveira@gmail.com}}

\begin{document}
\maketitle

\begin{abstract}
Safety cases for autonomous drones are typically static documents,
disconnected from the runtime behaviour they are meant to assure.
We present SCC (Safety Case Compiler), a tool that takes safety
requirements in a small DSL and compiles them into: (i) executable
Rust test stubs, (ii) Goal Structuring Notation (GSN) nodes, and
(iii) cryptographic proofs in a hash chain. One source of truth
(the requirements file) generates three verified artefacts. We
demonstrate SCC on a 3-requirement UAV envelope (speed, altitude,
obstacle clearance), showing that a single-line change regenerates
all artefacts and re-validates the chain. SCC is implemented in
Rust (\texttt{no\_std} compatible) and evaluated on 10 unit tests
and 1 end-to-end demo. We discuss how SCC integrates with the
DO-326A security process and identify open problems.
\end{abstract}

\begin{IEEEkeywords}
safety case, certification, GSN, Rust, UAV, DO-326A
\end{IEEEkeywords}

\section{Introduction}
% 1. Problem: safety cases are static, disconnected from runtime
% 2. Gap: no tooling from requirements to executable evidence
% 3. Contribution: SCC compiles requirements to 3 artefacts
% 4. Results: 10 tests, 1 demo, 3 artefacts per requirement

\section{Background}
\subsection{Safety Cases and GSN}
% GSN (Kelly 1998), static safety cases, EASA SC-VTOL
\subsection{Runtime Verification}
% Runtime monitors, CertChain (hash chain for safety events)
\subsection{Existing Tools}
% SafeCert, BusbySim — static; no compilation to executable proof

\section{SCC Design}
\subsection{DSL}
% requirements.safety: speed_max: 10.0 HIGH
\subsection{Compilation Pipeline}
% parse -> test stubs + GSN + chain blocks
\subsection{Why Rust}
% no_std, deterministic, hash chain (SHA-256)

\section{Implementation}
\subsection{Crate Structure}
% rnfc-scc (160 lines), 10 tests
\subsection{Integration with CertChain}
% Each requirement = one chain block

\section{Evaluation}
\subsection{Case Study: 3-Requirement Envelope}
% speed_max=10, altitude_range=5-120, obstacle_clearance=2
\subsection{Results}
% 3 requirements -> 3 tests + 3 GSN nodes + 3 chain blocks
% Single-line change regenerates all 3 artefacts
\subsection{Limitations}
% Only 3 requirement types; no quantifier support; synthetic physics

\section{Discussion}
\subsection{Relation to DO-326A}
% SCC provides evidence for SEC-* requirements
\subsection{Open Problems}
% Formal semantics for DSL; composition of safety cases

\section{Conclusion}
% SCC compiles requirements to executable evidence. Future:
% integrate with runtime monitors, support quantifiers.

\bibliographystyle{IEEEtran}
\bibliography{references}

\end{document}
